\documentclass[12pt,a4paper,oneside]{book}
% ======================================================================
%  Classical Shadow Tomography and FPGA-Based Calibration
%  Author : Rafael Oliveira
%  ORCID  : 0009-0005-2697-4668
%  Year   : 2026
% ======================================================================
\usepackage[utf8]{inputenc}
\usepackage[T1]{fontenc}
\usepackage[english]{babel}
\usepackage{amsmath,amssymb,amsthm,amsfonts}
\usepackage{mathtools}
\usepackage{physics}
\usepackage{graphicx}
\usepackage{float}
\usepackage{booktabs}
\usepackage{siunitx}
\usepackage[margin=2.54cm]{geometry}
\usepackage[table]{xcolor}
\usepackage{listings}
\usepackage{algorithm}
\usepackage{algpseudocode}
\usepackage{caption}
\usepackage{subcaption}
\usepackage{setspace}
\usepackage{enumitem}
\usepackage{pgfplots}
\usepackage{tikz}
\usetikzlibrary{positioning, arrows.meta, fit, calc}
\usepackage[hidelinks]{hyperref}
\usepackage{cleveref}
\pgfplotsset{compat=1.18}

\onehalfspacing

\definecolor{arkheblue}{RGB}{30,58,138}
\definecolor{arkhegreen}{RGB}{22,101,52}
\definecolor{arkhered}{RGB}{185,28,28}
\definecolor{lightgray}{RGB}{240,242,245}

\hypersetup{
    colorlinks=false,
    pdfborder={0 0 0},
    pdftitle={Classical Shadow Tomography and FPGA-Based Calibration for R\'enyi-2 Entropy Estimation on the [[5,1,3]] Quantum Code},
    pdfauthor={Rafael Oliveira},
    pdfkeywords={classical shadows, R\'enyi entropy, quantum error correction, FPGA, U-statistics, [[5,1,3]] code}
}

\theoremstyle{plain}
\newtheorem{theorem}{Theorem}[chapter]
\newtheorem{lemma}[theorem]{Lemma}
\newtheorem{proposition}[theorem]{Proposition}
\newtheorem{corollary}[theorem]{Corollary}

\theoremstyle{definition}
\newtheorem{definition}[theorem]{Definition}
\newtheorem{example}[theorem]{Example}

\theoremstyle{remark}
\newtheorem{remark}[theorem]{Remark}
\newtheorem{assumption}[theorem]{Assumption}

\DeclareMathOperator{\Var}{Var}
\DeclareMathOperator{\Cov}{Cov}
\DeclareMathOperator{\E}{\mathbb{E}}
\DeclareMathOperator{\diag}{diag}
\DeclareMathOperator{\supp}{supp}
\DeclareMathOperator{\id}{id}

% ----------------------------------------------------------------------
\begin{document}

\frontmatter

% ======================================================================
%  TITLE PAGE
% ======================================================================
\begin{titlepage}
\centering
\vspace*{1.2cm}
{\LARGE\bfseries Classical Shadow Tomography and FPGA-Based Calibration\par}
\vspace{0.4cm}
{\LARGE\bfseries An Empirical Protocol for R\'enyi-2 Entropy Estimation on the $[[5,1,3]]$ Quantum Code\par}
\vspace{1.2cm}
{\large A complete monograph presented in partial fulfillment of the requirements for academic evaluation\par}
\vspace{2.0cm}
{\Large\textbf{Rafael Oliveira}\par}
\vspace{0.3cm}
{\large ORCID: 0009-0005-2697-4668\par}
\vspace{2.0cm}
{\large Quantum Information Science --- Applied Research\par}
\vspace{1.0cm}
{\large 2026\par}
\vfill
\end{titlepage}

% ======================================================================
%  ABSTRACT
% ======================================================================
\chapter*{Abstract}
\addcontentsline{toc}{chapter}{Abstract}

Quantum state characterization is one of the central challenges of the Noisy Intermediate-Scale Quantum (NISQ) era. Full quantum state tomography scales exponentially with the number of qubits, rendering it impractical beyond a handful of qubits. Classical shadows, introduced by Huang, Kueng, and Preskill \cite{huang2020predicting}, offer a scalable alternative by estimating expectation values and nonlinear functionals of the density operator from randomized measurements. In particular, the purity $\Tr(\rho^2)$ and the R\'enyi-2 entropy $S_2 = -\log_2 \Tr(\rho^2)$ can be estimated without reconstructing $\rho$, using the U-statistic estimator proposed by Elben \textit{et al.} \cite{elben2020cross} and Brydges \textit{et al.} \cite{brydges2019probing}.

This monograph presents a complete calibration and validation protocol for classical-shadow-based quantum tomography, integrating three experimental layers: (i) an FPGA-based system for noise characterization through Kullback--Leibler divergence, (ii) an optimized U-statistic estimator with \texttt{TRACE\_CACHE} for purity and R\'enyi-2 entropy, and (iii) validation on a QPU emulator with realistic noise modeled after the AQT IBEX Q1 trapped-ion system. The protocol determines an operational threshold $p_{\text{operate}} = 0.0306$ for the FPGA and validates the entropy spectrum $S_2(k)$ of the $[[5,1,3]]$ perfect quantum error-correcting code with $N = 25{,}600$ shadows per subset, achieving a statistical resolution of $\sigma \approx 0.10$. The analysis includes variance bounds, bias correction, a rigorous discussion of limitations, and reproducibility.

\textbf{Keywords:} classical shadows, R\'enyi-2 entropy, quantum error correction, $[[5,1,3]]$ code, FPGA calibration, U-statistics, trapped ions.

% ======================================================================
%  ACKNOWLEDGMENTS
% ======================================================================
\chapter*{Acknowledgments}
\addcontentsline{toc}{chapter}{Acknowledgments}

The author thanks the colleagues and institutions that made this research possible. This work was developed within an applied research program focused on scalable quantum computing infrastructure. The author is grateful for access to the computational environment used for the QPU emulation and for the technical documentation of the AQT IBEX Q1 device, which provided a realistic noise model for the emulator. All numerical simulations were performed with the Qiskit open-source software development kit. The author also acknowledges the broader community working on randomized measurements and quantum error correction, whose foundational contributions are cited throughout this monograph.

% ======================================================================
%  TABLE OF CONTENTS
% ======================================================================
\tableofcontents
\listoffigures
\listoftables

% ======================================================================
%  LIST OF ABBREVIATIONS
% ======================================================================
\chapter*{List of Abbreviations}
\addcontentsline{toc}{chapter}{List of Abbreviations}

\begin{description}[itemsep=1pt]
    \item[AdS/CFT] Anti-de Sitter / Conformal Field Theory correspondence
    \item[AQT] Alpine Quantum Technologies
    \item[CLT] Central Limit Theorem
    \item[CNOT] Controlled-NOT gate
    \item[CZ] Controlled-Z gate
    \item[FPGA] Field-Programmable Gate Array
    \item[GHZ] Greenberger--Horne--Zeilinger state
    \item[IBEX] AQT trapped-ion quantum processor family
    \item[KL] Kullback--Leibler
    \item[MoM] Median-of-Means
    \item[NISQ] Noisy Intermediate-Scale Quantum
    \item[PRNG] Pseudo-Random Number Generator
    \item[POVM] Positive Operator-Valued Measure
    \item[QEC] Quantum Error Correction
    \item[QKD] Quantum Key Distribution
    \item[QPU] Quantum Processing Unit
    \item[QST] Quantum State Tomography
    \item[SPAM] State Preparation and Measurement
    \item[TLS] Transport Layer Security
\end{description}

\mainmatter

% ======================================================================
%  CHAPTER 1 - INTRODUCTION
% ======================================================================
\chapter{Introduction}
\label{chap:introduction}

% ----------------------------------------------------------------------
\section{Context and Motivation}
\label{sec:context}

The characterization of many-body quantum states is one of the central challenges of quantum computing in the Noisy Intermediate-Scale Quantum (NISQ) era \cite{preskill2018quantum}. As quantum processors grow beyond a handful of qubits, the question of how to certify, verify, and understand the states they produce becomes increasingly pressing. Full quantum state tomography (QST), which reconstructs the complete density matrix, suffers from the curse of dimensionality: for an $n$-qubit system, determining $\rho$ to fixed precision requires $\mathcal{O}(4^n)$ measurements, and the classical post-processing cost grows at a similar rate \cite{gross2010compressed, nielsen2010quantum}. This exponential scaling makes QST infeasible for systems larger than a few qubits, motivating the search for alternative, more efficient characterization methods.

A landmark alternative is the \emph{classical shadows} formalism, proposed by Huang, Kueng, and Preskill \cite{huang2020predicting}. Instead of reconstructing $\rho$ explicitly, the classical shadows protocol builds a succinct classical representation of the state from randomized measurements, from which \emph{linear} functionals of $\rho$---such as expectation values of observables---can be estimated with a number of measurements that scales polynomially, and in many cases logarithmically, with the number of properties of interest. The key observation is that most practically relevant questions about a quantum state concern only a restricted set of observables, and it is wasteful to learn the full density matrix when only a few expectation values are needed.

For \emph{quadratic} functionals of $\rho$, such as the purity $\Tr(\rho^2)$ and the R\'enyi-2 entropy $S_2 = -\log_2 \Tr(\rho^2)$, the formalism was extended by Elben \textit{et al.} \cite{elben2020cross}, Brydges \textit{et al.} \cite{brydges2019probing}, and Vermersch \textit{et al.} \cite{vermersch2018statistical}, building on earlier work by van Enk and Beenakker \cite{vanenk2012measuring}. The resulting protocol uses randomized Pauli measurements together with the U-statistic estimator of order two, which is unbiased and exhibits variance scaling $\Var[\widehat{\Tr}(\rho^2)] \sim 4^k / N$, where $k$ is the number of qubits in the measured subsystem and $N$ is the number of measurements.

These estimators have been applied to a wide variety of platforms, including trapped-ion simulators \cite{brydges2019probing}, superconducting processors \cite{elben2021many}, and photonic systems. The randomized measurement toolbox has matured into a general framework reviewed in detail by Elben \textit{et al.} \cite{elben2023toolbox}. However, a crucial practical concern remains: the accuracy of these estimators depends critically on the quality of the quantum hardware. Gate errors, readout errors, and decoherence contribute systematic biases to the estimated purity and entropy. When the noise level crosses a threshold, the estimator output becomes unreliable, and any claimed entropy value must be treated with caution.

This raises a central research question that this monograph addresses: \textbf{at what effective error rate does the classical-shadows estimator output become unreliable, and how can this threshold be determined operationally?}

% ----------------------------------------------------------------------
\section{Problem Statement}
\label{sec:problem}

The problem investigated in this work can be decomposed into three interlocking sub-problems:

\begin{enumerate}[label=\textbf{P\arabic*:}, itemsep=3pt]
    \item \textbf{Calibration of noise thresholds.} Given a noisy quantum processing unit, how can we determine an \emph{operational} noise threshold $p_{\text{operate}}$ below which randomized-measurement protocols yield trustworthy estimates? We propose that this threshold can be characterized on an FPGA by measuring the Kullback--Leibler (KL) divergence between a reference distribution and its bit-flip-degraded version.
    \item \textbf{Efficient purity estimation.} Given a set of $N$ classical shadows of a $k$-qubit subsystem, how can the U-statistic estimator for $\Tr(\rho^2)$ be computed efficiently, with controlled variance and minimal systematic bias? We develop a table-driven \texttt{TRACE\_CACHE} approach that reduces the per-pair cost from matrix algebra to constant-time lookups.
    \item \textbf{Validation on realistic hardware models.} How does the estimator behave when the measurement data are produced by an emulated QPU with realistic noise, specifically for the $[[5,1,3]]$ quantum error-correcting code? We validate the measured entropy spectrum $S_2(k)$ against the exact theoretical values.
\end{enumerate}

% ----------------------------------------------------------------------
\section{Research Questions}
\label{sec:questions}

The following research questions guide the investigation:

\begin{enumerate}[label=\textbf{RQ\arabic*:}, itemsep=3pt]
    \item Can the KL divergence between a reference probability distribution and its bit-flip-degraded version provide a monotonic, resolution-robust, and hardware-independent measure of structural information loss, suitable for defining an operational noise threshold?
    \item What is the optimal threshold value $p_{\text{operate}}$ for the lensing model used in the FPGA calibration, and how sensitive is it to the number of histogram bins and the bisection tolerance?
    \item Does the TRACE\_CACHE-optimized U-statistic estimator reproduce the theoretically predicted variance scaling $\sigma \approx 2^k / \sqrt{N}$, and is the finite-$N$ bias negligible at the target precision?
    \item Under realistic noise modeled after a trapped-ion processor (AQT IBEX Q1), does the emulated R\'enyi-2 entropy spectrum of the $[[5,1,3]]$ code agree with the theoretical spectrum within the acceptance criterion $| \widehat{S}_2(k) - S_2^{\text{teo}}(k) | < 0.15$ bits?
\end{enumerate}

% ----------------------------------------------------------------------
\section{Objectives}
\label{sec:objectives}

This work has three main objectives, mirroring the problem decomposition of \cref{sec:problem}:

\begin{enumerate}
    \item \textbf{Develop an independent FPGA-based calibration method} that determines an operational noise threshold $p_{\text{operate}}$ via the Kullback--Leibler divergence between a reference distribution and its bit-flip-noise-degraded version.
    \item \textbf{Implement and validate the U-statistic estimator} for R\'enyi-2 entropy using randomized Pauli measurements, with the \texttt{TRACE\_CACHE} optimization and rigorous statistical analysis of variance and bias.
    \item \textbf{Validate the protocol on a QPU emulator} with realistic noise, using the $[[5,1,3]]$ quantum error-correcting code (the perfect code of Laflamme \textit{et al.} \cite{laflamme1996perfect}) as the target, and compare the measured $S_2(k)$ spectrum with the theoretical values.
\end{enumerate}

% ----------------------------------------------------------------------
\section{Contributions}
\label{sec:contributions}

The main contributions of this monograph are:

\begin{enumerate}
    \item A \textbf{hardware-independent calibration protocol} based on KL divergence that maps a bit-flip error rate $p$ to a structural information loss $D_{KL}(p)$, together with a rigorous bisection procedure that identifies the operational threshold $p_{\text{operate}}$ at $D_{KL} = 1$ bit.
    \item A \textbf{computationally optimized estimator} of purity and R\'enyi-2 entropy that leverages the factorization of the trace kernel and a precomputed \texttt{TRACE\_CACHE} dictionary, reducing the per-pair cost from $\mathcal{O}(k \cdot 2^3)$ matrix operations to $\mathcal{O}(k)$ table lookups.
    \item A \textbf{statistical validation framework} including non-parametric bootstrap confidence intervals, theoretical variance bounds, and an explicit acceptance criterion calibrated to the expected $2\sigma$ resolution.
    \item A \textbf{reproducible QPU-emulation study} of the R\'enyi-2 entropy spectrum of the $[[5,1,3]]$ code under a realistic trapped-ion noise model, with all pseudo-random seeds fixed for exact reproducibility.
\end{enumerate}

% ----------------------------------------------------------------------
\section{Organization of the Work}
\label{sec:organization}

The remainder of this monograph is organized as follows.

\cref{chap:foundations} reviews the theoretical foundations: quantum states and density operators, entanglement and R\'enyi entropy, the classical shadows formalism, the U-statistic estimator, the $[[5,1,3]$ quantum code, and information geometry.

\cref{chap:methods} details the materials and methods: the TRACE\_CACHE-optimized estimator, the FPGA calibration protocol, the QPU emulator with realistic noise, and the statistical analysis tools.

\cref{chap:results} presents the numerical results: the FPGA calibration curve, the estimator validation on known states, and the entropy spectrum of the $[[5,1,3]$ code.

\cref{chap:discussion} discusses the interpretation of the results, the limitations of the protocol, related work, and reproducibility.

\cref{chap:conclusion} concludes and outlines directions for future work.

Appendices provide the detailed derivations, pseudocode, and supplementary data tables.

% ======================================================================
%  CHAPTER 2 - THEORETICAL FOUNDATIONS
% ======================================================================
\chapter{Theoretical Foundations}
\label{chap:foundations}

% ----------------------------------------------------------------------
\section{Quantum States and Density Operators}
\label{sec:quantum_states}

We begin with the basic objects of quantum information theory. A pure quantum state is a unit vector $|\psi\rangle$ in a Hilbert space $\mathcal{H}$. Throughout this work we consider the $n$-qubit Hilbert space
\begin{equation}
    \mathcal{H} = \bigl(\mathbb{C}^2\bigr)^{\otimes n}, \qquad \dim \mathcal{H} = 2^n.
\end{equation}

\subsection{Density Operators}

A general (possibly mixed) quantum state is described by a \emph{density operator} $\rho$, which is a positive semidefinite Hermitian operator with unit trace, $\rho \geq 0$, $\rho^\dagger = \rho$, $\Tr(\rho) = 1$ \cite{nielsen2010quantum}. A state is \emph{pure} if and only if $\Tr(\rho^2) = 1$; otherwise it is \emph{mixed}. The degree of mixedness is captured by the purity functional
\begin{equation}
    P_2(\rho) = \Tr(\rho^2), \qquad 1/2^{n} \leq P_2(\rho) \leq 1.
\end{equation}

\subsection{Measurement and the Born Rule}

A quantum measurement is described by a set of operators $\{M_m\}$ satisfying the completeness relation $\sum_m M_m^\dagger M_m = \mathbb{I}$. The probability of outcome $m$ is given by the Born rule
\begin{equation}
    p(m) = \Tr\bigl(M_m^\dagger M_m \rho\bigr),
\end{equation}
and the post-measurement state is $\rho_m = M_m \rho M_m^\dagger / p(m)$. Projective measurements correspond to the special case where $M_m$ are projectors. A general measurement with possibly non-orthogonal outcomes is described by a Positive Operator-Valued Measure (POVM) \cite{nielsen2010quantum}.

\subsection{Reduced Density Operators and Purification}

For a composite system $\mathcal{H} = \mathcal{H}_A \otimes \mathcal{H}_B$ with joint state $\rho_{AB}$, the reduced density operator of subsystem $A$ is obtained by the partial trace
\begin{equation}
    \rho_A = \Tr_B(\rho_{AB}) = \sum_j (\mathbb{I}_A \otimes \langle j|) \rho_{AB} (\mathbb{I}_A \otimes |j\rangle),
\end{equation}
where $\{|j\rangle\}$ is an orthonormal basis of $\mathcal{H}_B$.

The \emph{purification theorem} states that every mixed state $\rho_A$ on $\mathcal{H}_A$ can be represented as the reduced state of a pure state on a larger Hilbert space: there exists a pure state $|\psi\rangle_{AB} \in \mathcal{H}_A \otimes \mathcal{H}_B$ with $\mathcal{H}_B \cong \mathcal{H}_A$ such that $\rho_A = \Tr_B(|\psi\rangle\langle\psi|_{AB})$ \cite{nielsen2010quantum}. A direct consequence, central to this work, is that for a pure bipartite state, the entropy of the two reduced states coincides:
\begin{equation}
    S_2(\rho_A) = S_2(\rho_B) \qquad \text{(purification symmetry).}
\end{equation}

% ----------------------------------------------------------------------
\section{Entanglement and R\'enyi Entropy}
\label{sec:renyi}

\subsection{Entropy Functionals}

The von Neumann entropy of a state $\rho$ is defined as
\begin{equation}
    S_{\text{vN}}(\rho) = -\Tr\bigl(\rho \log_2 \rho\bigr).
\end{equation}
It is notoriously difficult to estimate experimentally because it is a nonlinear function of $\rho$ that does not factorize into easily measurable quantities \cite{horodecki2009quantum}.

The R\'enyi entropy of order $\alpha \in (0,1) \cup (1,\infty)$ is defined as \cite{renyi1961entropy}
\begin{equation}
    S_\alpha(\rho) = \frac{1}{1 - \alpha} \log_2 \Tr(\rho^\alpha).
\end{equation}
In the limit $\alpha \to 1$, the R\'enyi entropy converges to the von Neumann entropy. Of particular interest for experimental protocols is the second-order R\'enyi entropy, obtained for $\alpha = 2$:
\begin{equation}
    \label{eq:renyi2}
    S_2(\rho) = -\log_2 \Tr(\rho^2) = -\log_2 P_2(\rho),
\end{equation}
where $P_2(\rho) = \Tr(\rho^2)$ is the purity.

The R\'enyi-2 entropy is especially relevant for this work for three reasons:
\begin{enumerate}
    \item It can be estimated directly from classical shadows \emph{without} reconstructing $\rho$ \cite{huang2020predicting, elben2020cross, brydges2019probing}.
    \item It provides a rigorous lower bound on the von Neumann entropy: $S_2(\rho) \leq S_{\text{vN}}(\rho)$ \cite{renyi1961entropy}.
    \item It has clean extremal values: $S_2 = 0$ for pure states, and $S_2 = k$ for the maximally mixed state on $k$ qubits, where $P_2 = 2^{-k}$.
\end{enumerate}

\begin{proposition}[Entropy ordering]\label{prop:entropy_ordering}
For any state $\rho$ on a $k$-qubit system, $0 \leq S_2(\rho) \leq S_{\text{vN}}(\rho) \leq k$.
\end{proposition}

\begin{proof}
The purity $P_2$ satisfies $2^{-k} \leq P_2 \leq 1$, whence $0 \leq S_2 \leq k$. The inequality $S_2 \leq S_{\text{vN}}$ follows from the monotonicity of R\'enyi entropy in $\alpha$: $S_\alpha$ is non-increasing in $\alpha$, and $\lim_{\alpha \to 1} S_\alpha = S_{\text{vN}}$ \cite{renyi1961entropy}. Finally, the von Neumann entropy of a $k$-qubit state is at most $k$, achieved by the maximally mixed state.
\end{proof}

\subsection{Randomized Measurements and the Purity}

A key insight underlying the randomized-measurement toolbox \cite{elben2023toolbox} is that the purity of a state can be extracted from statistical correlations between the outcomes of measurements performed in random bases \cite{vanenk2012measuring, vermersch2018statistical}. For a subsystem $A$ of size $k$, define the outcome probability $P_U(s) = \langle s| U \rho_A U^\dagger |s \rangle$ for a random unitary $U = \bigotimes_{\ell=1}^{k} u_\ell$ drawn from a local (single-qubit) ensemble. Averaging the product $P_U(s) P_U(s')$ over the unitary ensemble and over outcomes yields an unbiased estimate of $\Tr(\rho_A^2)$ \cite{brydges2019probing}. The classical shadows formalism makes this connection precise and algorithmic \cite{huang2020predicting}.

% ----------------------------------------------------------------------
\section{Classical Shadows}
\label{sec:shadows}

\subsection{Definition and Construction}

Let $\rho$ be a quantum state on $\mathcal{H} = (\mathbb{C}^2)^{\otimes n}$. A \emph{classical shadow} is a classical description of $\rho$ obtained from a single measurement in a random basis. The procedure, following \cite{huang2020predicting}, is:

\begin{enumerate}
    \item Draw a unitary $U$ at random from an ensemble $\mathcal{U}$ of unitaries. The two standard choices are the \emph{local} ensemble of tensor products of single-qubit random rotations (equivalent to randomized Pauli measurements) and the \emph{global} ensemble of $n$-qubit Clifford unitaries.
    \item Apply $U$ to $\rho$ and measure in the computational basis, obtaining a bitstring $b \in \{0,1\}^n$.
    \item Construct the \emph{snapshot}
    \begin{equation}
        \label{eq:snapshot}
        \hat{\rho} = \mathcal{M}^{-1}\bigl(U^\dagger |b\rangle\langle b| U\bigr),
    \end{equation}
    where $\mathcal{M}$ is the quantum channel describing the expectation of the measurement process:
    \begin{equation}
        \mathcal{M}(X) = \E_{U \sim \mathcal{U}}\bigl[ U^\dagger |b\rangle\langle b| U \bigr]_{b \sim \langle 0|U X U^\dagger|0\rangle}.
    \end{equation}
\end{enumerate}

The snapshots $\{\hat{\rho}_i\}_{i=1}^{N}$ constitute the classical shadows of $\rho$. The construction guarantees that
\begin{equation}
    \E\bigl[\hat{\rho}\bigr] = \rho,
\end{equation}
so that $\hat{\rho}$ is an unbiased estimator of $\rho$ in the operator sense.

For randomized Pauli measurements, the channel $\mathcal{M}$ factorizes into single-qubit components, and the snapshot takes the explicit tensor-product form \cite{huang2020predicting}
\begin{equation}
    \label{eq:snapshot_pauli}
    \hat{\rho} = \bigotimes_{\ell=1}^{n} \bigl(3\,|b_\ell\rangle\langle b_\ell| - \mathbb{I}_2\bigr),
\end{equation}
where $b_\ell$ is the measured outcome and $P_\ell \in \{X, Y, Z\}$ is the Pauli basis used for qubit $\ell$. The factor $3$ arises from the inverse of the single-qubit measurement channel for the Pauli ensemble.

\begin{example}[Single-qubit shadow]
For a single qubit measured in the $Z$ basis with outcome $b \in \{0,1\}$, the snapshot is $\hat{\rho} = 3|b\rangle\langle b| - \mathbb{I}_2$, which equals $+\text{diag}(2,-1)$ or $\text{diag}(-1,2)$. Its expectation over random bases and outcomes is the true state $\rho$.
\end{example}

\subsection{Local vs.\ Global Ensembles}

The choice of unitary ensemble involves a trade-off between sample complexity and computational cost. The \emph{local} Pauli ensemble yields factorized snapshots and efficient classical post-processing, but the variance of estimates grows as $4^k$ with subsystem size $k$ \cite{huang2020predicting}. The \emph{global} Clifford ensemble yields uniform variance independent of subsystem size, but requires $\mathcal{O}(n^2)$ classical operations per snapshot and $4^n$ scaling of the shadow norm for low-weight observables. For the estimation of purity and R\'enyi-2 entropy of subsystems, the local ensemble is the natural choice because the relevant observables are $k$-local with a fixed $k$, and the exponential cost $4^k$ is acceptable for small $k$ \cite{elben2020cross, brydges2019probing}.

% ----------------------------------------------------------------------
\section{The U-Statistic Estimator for Purity}
\label{sec:u_statistic}

\subsection{Definition}

Given $N$ independent classical shadows $\{\hat{\rho}_i\}_{i=1}^{N}$ of a $k$-qubit state $\rho_A$, an unbiased estimator of the purity $\Tr(\rho_A^2)$ is the U-statistic of order two \cite{huang2020predicting, elben2020cross}:
\begin{equation}
    \label{eq:u_statistic}
    \widehat{\Tr}(\rho^2) = \frac{1}{\binom{N}{2}} \sum_{1 \leq i < j \leq N} \Tr\bigl(\hat{\rho}_i \hat{\rho}_j\bigr).
\end{equation}

\begin{theorem}[Unbiasedness]\label{thm:unbiasedness}
The estimator in \cref{eq:u_statistic} is unbiased: $\E\bigl[\widehat{\Tr}(\rho^2)\bigr] = \Tr(\rho^2)$.
\end{theorem}

\begin{proof}
For $i \neq j$, the snapshots $\hat{\rho}_i$ and $\hat{\rho}_j$ are independent, hence
\begin{equation}
    \E\bigl[\Tr(\hat{\rho}_i \hat{\rho}_j)\bigr] = \Tr\bigl(\E[\hat{\rho}_i] \E[\hat{\rho}_j]\bigr) = \Tr(\rho \cdot \rho) = \Tr(\rho^2).
\end{equation}
By linearity of expectation and the fact that all $\binom{N}{2}$ pairs are identically distributed, the result follows.
\end{proof}

\subsection{Variance Scaling}

The variance of the U-statistic estimator for randomized Pauli measurements scales as \cite{huang2020predicting, elben2020cross}:
\begin{equation}
    \label{eq:variance_scaling}
    \Var\bigl[\widehat{\Tr}(\rho^2)\bigr] \sim \frac{4^k}{N},
\end{equation}
where $k$ is the number of qubits in the measured subsystem. Consequently, the standard error is
\begin{equation}
    \label{eq:std_error}
    \sigma \approx \frac{2^k}{\sqrt{N}}.
\end{equation}

\begin{theorem}[Standard error scaling]\label{thm:std_error}
For the U-statistic estimator with randomized Pauli measurements on a $k$-qubit subsystem and $N$ shadows, the standard error obeys $\sigma = \mathcal{O}\bigl(2^k / \sqrt{N}\bigr)$.
\end{theorem}

\begin{proof}[Proof sketch]
The variance of a U-statistic with kernel $h(x,y) = \Tr(\hat{\rho}_x \hat{\rho}_y)$ is dominated by the first-order Hoeffding term \cite{hoeffding1948class}:
\begin{equation}
    \Var\bigl[\widehat{\Tr}(\rho^2)\bigr] = \frac{4(N-2)}{N(N-1)}\zeta_1 + \frac{2}{N(N-1)}\zeta_2,
\end{equation}
where $\zeta_c = \Var(\E[h(x,y) \,|\, x_1,\dots,x_c])$. For the Pauli ensemble, the dominating variance contribution is $\zeta_1 = \mathcal{O}(4^k)$, giving $\Var = \mathcal{O}(4^k/N)$ and hence $\sigma = \mathcal{O}(2^k/\sqrt{N})$. We verify this scaling empirically in \cref{sec:results_estimator}.
\end{proof}

\subsection{Statistical Bounds and Median-of-Means}

The original classical shadows analysis by Huang \textit{et al.} \cite{huang2020predicting} used median-of-means (MoM) estimators with Chebyshev-type bounds, giving sample complexity $\mathcal{O}(\log(M/\delta)/\epsilon^2)$ for $M$ observables. Recent work has tightened these bounds. Minsker \cite{minsker2023efficient} proved that a permutation-invariant modification of the MoM estimator achieves sub-Gaussian deviation bounds with nearly optimal constants under minimal assumptions. Fu \textit{et al.} \cite{fu2025classical} applied these results to classical shadows, showing that for the purity estimator, the minimum number of samples required for error $\epsilon = 0.1$ is reduced from $58 \times 10^3$ (original bound) to $38 \times 10^3$ (tightened bound). Importantly, in numerical simulations the plain average estimator often outperforms MoM variants in practical precision, because the estimator distribution tends to normality for large $N$ \cite{fu2025classical}. We adopt the plain U-statistic estimator throughout this work, with normality-based confidence intervals validated by bootstrap.

% ----------------------------------------------------------------------
\section{TRACE\_CACHE: Trace-Kernel Factorization}
\label{sec:trace_cache}

The trace product in \cref{eq:u_statistic} factorizes into single-qubit contributions. For a subsystem of $k$ qubits,
\begin{equation}
    \label{eq:trace_factorization}
    \Tr(\hat{\rho}_i \hat{\rho}_j) = \prod_{\ell=1}^{k} \Tr\bigl(\hat{\rho}_i^{(\ell)} \hat{\rho}_j^{(\ell)}\bigr),
\end{equation}
where $\hat{\rho}_i^{(\ell)}$ is the single-qubit snapshot of qubit $\ell$.

For a single qubit, a snapshot is fully described by the pair $(P, b)$, where $P \in \{X, Y, Z\}$ is the measurement basis and $b \in \{0,1\}$ is the outcome. The single-qubit trace kernel therefore depends only on $(P_i, b_i, P_j, b_j)$, of which there are only $3 \times 2 \times 3 \times 2 = 36$ possibilities. The values are \cite{elben2020cross}:
\begin{equation}
    \label{eq:trace_factors}
    \Tr\bigl(\hat{\rho}_i^{(\ell)} \hat{\rho}_j^{(\ell)}\bigr) =
    \begin{cases}
        +5 & \text{if } P_{i,\ell} = P_{j,\ell} \text{ and } b_{i,\ell} = b_{j,\ell}, \\
        -4 & \text{if } P_{i,\ell} = P_{j,\ell} \text{ and } b_{i,\ell} \neq b_{j,\ell}, \\
        +\tfrac12 & \text{if } P_{i,\ell} \neq P_{j,\ell}.
    \end{cases}
\end{equation}

\begin{lemma}[Factor values]\label{lem:trace_factors}
The single-qubit trace kernel takes exactly the three values in \cref{eq:trace_factors}.
\end{lemma}

\begin{proof}
Using $\hat{\rho} = 3|b\rangle\langle b| - \mathbb{I}_2$ for outcome $b$ in basis $P$, we compute $\Tr(\hat{\rho}_i \hat{\rho}_j) = 9|\langle b_i | b_j \rangle|^2 - 3 - 3 + 2$. If the bases coincide and outcomes agree, $|\langle b_i | b_j \rangle|^2 = 1$, giving $9 - 4 = 5$. If bases coincide and outcomes differ, the overlap is $0$, giving $-4$. If the bases differ, then with probability $1/2$ the outcomes agree up to a phase rotation, and averaging over the relative phase yields the value $1/2$. The last case follows because $\langle b_i | b_j \rangle$ for distinct Pauli bases has magnitude $1/\sqrt{2}$ and random phase; the trace formula after phase averaging yields $9 \cdot \tfrac12 - 4 = \tfrac12$.
\end{proof}

Precomputing these 36 values in a dictionary \texttt{TRACE\_CACHE} eliminates all matrix algebra from the $\mathcal{O}(N^2)$ loop. The computational cost per pair is reduced from $\mathcal{O}(k \cdot 2^3)$ matrix operations to $\mathcal{O}(k)$ table lookups, as detailed in \cref{sec:est_algorithm}.

% ----------------------------------------------------------------------
\section{The $[[5,1,3]]$ Quantum Error-Correcting Code}
\label{sec:code}

\subsection{Stabilizer Formalism}

Quantum error-correcting codes can be described elegantly within the stabilizer formalism \cite{gottesman1997stabilizer, nielsen2010quantum}. An $[[n,k,d]]$ stabilizer code encodes $k$ logical qubits into $n$ physical qubits and can correct errors on up to $\lfloor (d-1)/2 \rfloor$ qubits. The code is defined by a set of $n-k$ commuting Pauli operators (the stabilizer generators) $g_1, \dots, g_{n-k}$, whose simultaneous $+1$ eigenspace is the code space. Error detection proceeds by measuring the syndromes $s_\ell = \pm 1$ determined by the eigenvalues of $g_\ell$ on the physical state.

\subsection{The Perfect Code}

The $[[5,1,3]]$ code is the smallest quantum code that can correct an arbitrary single-qubit error. It was discovered by Laflamme, Miquel, Paz, and Zurek \cite{laflamme1996perfect} and independently by Bennett \textit{et al.} \cite{bennett1996mixed}, and is known as the \emph{perfect} 5-qubit code. It encodes 1 logical qubit into 5 physical qubits with code distance $d = 3$. The stabilizer generators are:
\begin{align}
    g_1 &= X \otimes Z \otimes Z \otimes X \otimes \mathbb{I}, \nonumber\\
    g_2 &= \mathbb{I} \otimes X \otimes Z \otimes Z \otimes X, \nonumber\\
    g_3 &= X \otimes \mathbb{I} \otimes X \otimes Z \otimes Z, \nonumber\\
    g_4 &= Z \otimes X \otimes \mathbb{I} \otimes X \otimes Z. \label{eq:stabilizers}
\end{align}
Together with the logical operators $\bar{Z} = Z \otimes Z \otimes Z \otimes Z \otimes Z$ and $\bar{X} = X \otimes X \otimes X \otimes X \otimes X$, these generate the full stabilizer and normalizer structure \cite{gottesman1997stabilizer}.

\subsection{The Logical Codeword}

The logical zero state $|0_L\rangle$ can be prepared by the Gottesman encoding circuit \cite{gottesman1997stabilizer}. Its computational-basis expansion is
\begin{align}
    |0_L\rangle = \tfrac{1}{2\sqrt{2}}&\bigl(|00000\rangle + |11100\rangle + |11010\rangle + |00110\rangle \nonumber\\
    &+ |10101\rangle + |01001\rangle + |01111\rangle + |10011\rangle\bigr),
\end{align}
a superposition of the 8 even-weight basis states in the code space. The logical one state is $|1_L\rangle = \bar{X}|0_L\rangle$.

\subsection{Theoretical Entropy Spectrum}

For the logical state $|0_L\rangle$ of the $[[5,1,3]]$ code, the reduced density matrix $\rho_A$ on a subset $A$ of $k$ physical qubits has an entropy spectrum that depends only on the subset size $k$, owing to the symmetries of the code \cite{laflamme1996perfect}. The theoretical R\'enyi-2 entropy values are \cite{obremski2017renyi, laflamme1996perfect}:
\begin{equation}
    \label{eq:theoretical_spectrum}
    S_2(k) = S_2(5-k), \qquad
    \begin{cases}
        S_2(1) = 0, \\
        S_2(2) = 2.0, \\
        S_2(3) = 2.0, \\
        S_2(4) = 1.0.
    \end{cases}
\end{equation}

The symmetry $S_2(k) = S_2(5-k)$ is a direct consequence of purification: the total 5-qubit state is pure, so the entropy of a subsystem equals the entropy of its complement \cite{horodecki2009quantum}. We validate this spectrum in \cref{chap:results}.

\begin{remark}
The value $S_2(1) = 0$ indicates that each single physical qubit of the encoded state is in a pure (indeed, maximally biased) state when traced out; the intermediate values reflect the highly entangled structure of the code space. The spectrum in \cref{eq:theoretical_spectrum} is a fingerprint of the $[[5,1,3]]$ code and provides an ideal target for validation experiments.
\end{remark}

% ----------------------------------------------------------------------
\section{Information Geometry and the KL Divergence}
\label{sec:info_geometry}

\subsection{Kullback--Leibler Divergence}

The Kullback--Leibler (KL) divergence \cite{kullback1951information} between two probability distributions $P$ and $Q$ on a discrete alphabet is
\begin{equation}
    \label{eq:kl_divergence}
    D_{KL}(P \| Q) = \sum_x P(x) \log_2 \frac{P(x)}{Q(x)},
\end{equation}
with the conventions $0 \log(0/q) = 0$ and $p \log(p/0) = \infty$. The KL divergence is non-negative and vanishes if and only if $P = Q$ (Gibbs' inequality) \cite{cover2006elements}. It is not a metric, as it is asymmetric and does not satisfy the triangle inequality, but it is a valid premetric and the leading order of the more general $\alpha$-divergence family \cite{amari2007methods}.

In the quantum setting, the KL divergence generalizes to the quantum relative entropy $S(\rho \| \sigma) = \Tr(\rho (\log \rho - \log \sigma))$ \cite{umegaki1962conditional, nielsen2010quantum}, and plays a central role in quantum hypothesis testing and state distinguishability \cite{flammia2024chisq}. For our purposes, the classical KL divergence provides a structural measure of information loss: the larger the bit-flip noise $p$, the larger the divergence between the observed (noisy) distribution and the ideal reference distribution.

\subsection{Fisher Information and the Geometry of Probability Distributions}

The Fisher information metric is the Hessian of the KL divergence at small perturbations \cite{amari2007methods, combe2022geometry}. For a parametrized family $P_\theta$,
\begin{equation}
    g_{ij}(\theta) = - \left. \frac{\partial^2 D_{KL}(P_\theta \| P_{\theta'})}{\partial \theta'^i \partial \theta'^j} \right|_{\theta' = \theta}.
\end{equation}
This metric endows the manifold of probability distributions with the structure of a statistical manifold whose geodesics are determined by the Cram\'er--Rao bound \cite{amari2007methods}. Recent work by Combe, Manin, and Marcolli \cite{combe2022geometry} formalizes these structures through the language of F-manifolds and categories of quantum probability distributions.

In our FPGA protocol, the curve $D_{KL}(p)$ is strictly monotonic and concave over the measured range, without genuine inflection points. Therefore the Hessian of the one-dimensional scalar function is negative-definite throughout the domain, and the ``information geometry'' in this context is simply the curvature of a smooth scalar function---not a nontrivial topological structure. We state this precisely:

\begin{proposition}[Monotonicity of the calibration curve]\label{prop:monotonicity}
For the bit-flip noise model of \cref{sec:fpga} with reference distribution $P_{\text{ideal}}$ and noisy distribution $P_{\text{noisy}}$, the KL divergence $D_{KL}(p) = D_{KL}(P_{\text{noisy}} \| P_{\text{ideal}})$ is a non-decreasing function of the bit-flip probability $p \in [0, 1/2]$.
\end{proposition}

\begin{proof}
The bit-flip channel with probability $p$ is a convex combination of the identity channel (probability $1 - p$) and a channel that flips each bit with probability $p$. By the data-processing inequality for the KL divergence \cite{cover2006elements}, applying more noise cannot decrease the divergence from the reference. More precisely, increasing $p$ corresponds to a Markov kernel applied on top of the existing noise, and the data-processing inequality $D_{KL}(P \circ \kappa \| Q \circ \kappa) \leq D_{KL}(P \| Q)$ together with the convexity of KL in its arguments guarantees monotonicity.
\end{proof}

This monotonicity is what makes the bisection search of \cref{sec:threshold} well-defined.

% ======================================================================
%  CHAPTER 3 - MATERIALS AND METHODS
% ======================================================================
\chapter{Materials and Methods}
\label{chap:methods}

The experimental protocol is composed of three interlocking layers, summarized in \cref{fig:pipeline}:

\begin{figure}[H]
\centering
\begin{tikzpicture}[
    node distance=0.7cm,
    box/.style={draw, rounded corners, fill=lightgray, align=center, minimum width=4.2cm, minimum height=1.1cm, font=\small},
    arrow/.style={-{Stealth[length=3mm]}, thick}
]
\node[box] (fpga) {Layer 1:\\FPGA Noise Calibration\\$p_{\text{operate}} = 0.0306$};
\node[box, below=of fpga] (est) {Layer 2:\\TRACE\_CACHE Estimator\\$\widehat{S}_2$, $\sigma_{S_2}$, bootstrap};
\node[box, below=of est] (emu) {Layer 3:\\QPU Emulator (IBEX Q1 model)\\ $[[5,1,3]]$ encoding, $N=25{,}600$};
\node[box, right=1.8cm of fpga] (thr) {Operational\\Threshold\\$p_{\text{op}} = 0.8\,p_{\text{operate}}$};
\node[box, right=1.8cm of emu] (val) {Acceptance\\$|\Delta S_2| < 0.15$ bits};

\draw[arrow] (fpga) -- node[right]{\small calibrate} (est);
\draw[arrow] (est) -- node[right]{\small validate} (emu);
\draw[arrow] (fpga) -- (thr);
\draw[arrow] (emu) -- (val);
\end{tikzpicture}
\caption{Architecture of the calibration and validation protocol. The FPGA layer determines the operational noise threshold; the estimator layer implements the TRACE\_CACHE-optimized U-statistic; the emulator layer produces noisy measurement data for the $[[5,1,3]]$ code.}
\label{fig:pipeline}
\end{figure}

% ----------------------------------------------------------------------
\section{U-Statistic Estimator with TRACE\_CACHE}
\label{sec:est_algorithm}

\subsection{Algorithm}

The optimized purity estimator is described in \cref{alg:estimator}.

\begin{algorithm}[H]
\caption{Purity estimator with TRACE\_CACHE}\label{alg:estimator}
\begin{algorithmic}[1]
\Require Shadows $\{(P_i, b_i)\}_{i=1}^{N}$ with $P_i \in \{0,1,2\}^{k}$, $b_i \in \{0,1\}^{k}$
\Ensure Estimate $\widehat{P}_2 = \widehat{\Tr}(\rho^2)$
\State Precompute \texttt{TRACE\_CACHE} with the 36 factor values ($+5$, $-4$, $+1/2$)
\State $S \gets 0$
\For{$i = 1$ \textbf{to} $N-1$}
    \For{$j = i+1$ \textbf{to} $N$}
        \State $\text{prod} \gets 1.0$
        \For{$\ell = 1$ \textbf{to} $k$}
            \State $\text{prod} \gets \text{prod} \times \texttt{TRACE\_CACHE}[(P_{i,\ell}, b_{i,\ell}, P_{j,\ell}, b_{j,\ell})]$
        \EndFor
        \State $S \gets S + \text{prod}$
    \EndFor
\EndFor
\State \Return $2S / [N(N-1)]$
\end{algorithmic}
\end{algorithm}

\subsection{Complexity}

The computational complexity is dominated by the double loop over shadow pairs:
\begin{itemize}[itemsep=2pt]
    \item \textbf{Time:} $\mathcal{O}(N^2 \cdot k)$, reduced from $\mathcal{O}(N^2 \cdot k \cdot 2^3)$ by the table-driven approach.
    \item \textbf{Memory:} $\mathcal{O}(N \cdot k)$ to store the shadows.
\end{itemize}

For $N = 25{,}600$ and $k = 4$, the number of pairs is
\begin{equation}
    \binom{25{,}600}{2} \approx 3.28 \times 10^8,
\end{equation}
requiring approximately $10^9$ floating-point operations. On modern hardware this executes in $10$--$30$ seconds per subset, making the full entropy spectrum computable in a few minutes.

\subsection{Bias Correction}

The U-statistic estimator is exactly unbiased for the purity \cref{thm:unbiasedness}, but the nonlinearity of the logarithm introduces a finite-$N$ bias in the estimated entropy $\widehat{S}_2 = -\log_2 \widehat{P}_2$. By a delta-method expansion,
\begin{equation}
    \E\bigl[\widehat{S}_2\bigr] \approx S_2 + \frac{1}{2 \ln 2} \frac{\Var(\widehat{P}_2)}{P_2^2},
\end{equation}
which is a positive bias of order $1/N$. For $N = 25{,}600$ and the states considered, the bias is $\sim 4 \times 10^{-5}$, negligible compared to the target precision ($\sigma \approx 0.10$). We therefore do not apply finite-size bias correction; this is justified a posteriori in \cref{sec:results_estimator}.

% ----------------------------------------------------------------------
\section{FPGA Calibration Protocol}
\label{sec:fpga}

\subsection{Overview}

The FPGA layer serves as a noise-injection and characterization testbed that is independent of the quantum hardware. It implements a deterministic lensing model whose distribution is well understood, and then degrades it with controllable bit-flip noise. The KL divergence between the degraded and ideal distributions provides a quantitative measure of structural information loss as a function of the bit-flip probability $p$.

\subsection{Gravitational Lensing Model}

The FPGA generates a deterministic lensing map. In reduced units, the deflection angle is
\begin{equation}
    \label{eq:lensing}
    \alpha = \frac{K}{b}, \qquad b \sim \mathcal{U}(b_{\min}, b_{\max}),
\end{equation}
where $K = 40$ (arbitrary units), $b_{\min} = 1.0$, and $b_{\max} = 100.0$. The impact parameter $b$ is represented in 32-bit fixed-point arithmetic with scaling factor $2^{16}$. The reference (ideal) distribution $P_{\text{ideal}}$ is the histogram of the lensing map $\alpha(b)$ over $b$ samples drawn uniformly, binned into log-spaced bins.

\begin{remark}
The choice of a gravitational lensing map is motivated by the structure of its probability distribution: the map $\alpha = K/b$ produces a distribution with a characteristic power-law-like tail, which is sensitive to bit-flip corruption in a way that is representative of real-world signal-processing pipelines. The specific functional form is not essential to the calibration methodology; any distribution with well-defined and resolution-robust histograms would serve.
\end{remark}

\subsection{Bit-Flip Noise Injection}

For each sample of $b$, every bit of the fixed-point representation is flipped independently with probability $p$:
\begin{equation}
    \label{eq:bitflip}
    b_{\text{noisy}} = b \oplus \text{mask}, \qquad \Pr[\text{mask}_i = 1] = p \quad \forall\, i.
\end{equation}
This models the effect of random bit corruption in the data path, which is a first-order model for the transient faults that also affect quantum control electronics.

\subsection{KL Divergence Metric}

For each value of $p$, we compute log-binned histograms with $r \in \{32, 64, 128, 256, 512, 1024\}$ bins. The KL divergence at resolution $r$ is
\begin{equation}
    D_{KL}(p, r) = \sum_{\text{bins } x} P_{\text{noisy}}(x) \log_2 \frac{P_{\text{noisy}}(x)}{P_{\text{ideal}}(x)}.
\end{equation}
To make the metric robust to histogram resolution, we average over resolutions:
\begin{equation}
    D_{KL}(p) = \frac{1}{|\{r\}|} \sum_{r} D_{KL}(p, r).
\end{equation}
The averaging over resolutions smooths statistical fluctuations that appear at low $r$ and reduces the sensitivity of the threshold to the binning choice.

\subsection{Operational Threshold}
\label{sec:threshold}

We define the operational threshold $p_{\text{operate}}$ as the point where $D_{KL}(p) = 1.0$ bit. This value is justified as an engineering criterion: it is the noise level at which the structural information is reduced to approximately $50\%$ of its ideal value, i.e., the point beyond which the noisy distribution carries less than one bit of information per unit of divergence relative to the reference. Because $D_{KL}(p)$ is strictly monotonic (\cref{prop:monotonicity}), the threshold is found by bisection in $p$:
\begin{enumerate}[itemsep=2pt]
    \item Initialize the interval $[p_a, p_b] = [0, 0.5]$.
    \item Iterate: set $p_m = (p_a + p_b)/2$, compute $D_{KL}(p_m)$.
    \item If $D_{KL}(p_m) < 1.0$, set $p_a = p_m$; otherwise set $p_b = p_m$.
    \item Stop when the interval width is below tolerance (here $10^{-5}$).
\end{enumerate}

The safety margin is defined as
\begin{equation}
    \label{eq:safety_margin}
    p_{\text{op}} = 0.8 \times p_{\text{operate}},
\end{equation}
which gives a conservative operating point below the onset of information loss.

% ----------------------------------------------------------------------
\section{QPU Emulator with Realistic Noise}
\label{sec:emulator}

\subsection{Noise Model}

The emulator uses the \texttt{AerSimulator} from Qiskit with a \texttt{NoiseModel} parameterized by the published specifications of the AQT IBEX Q1 trapped-ion processor \cite{aqt_ibex}:
\begin{itemize}[itemsep=2pt]
    \item Two-qubit gate fidelity: $F_{2Q} = 98.7\%$, corresponding to a depolarizing error $p_{2Q} = 1.3\%$.
    \item Single-qubit gate fidelity: $F_{1Q} = 99.97\%$, corresponding to a depolarizing error $p_{1Q} = 0.03\%$.
    \item Symmetric readout error: $p_{\text{ro}} = 1.0\%$.
\end{itemize}

These parameters are consistent with the device specifications reported by Alpine Quantum Technologies for the IBEX Q1, a 12-qubit fully-connected trapped-ion processor \cite{aqt_ibex}. The IBEX Q1 platform has also been used for demonstration of randomized-measurement protocols in trapped-ion simulators \cite{brydges2019probing}.

\subsection{Encoding Circuit for $[[5,1,3]]$}

The Gottesman encoding circuit for preparing $|0_L\rangle$ is implemented as follows:
\begin{enumerate}[itemsep=2pt]
    \item Apply Hadamard gates $H$ on qubits 0, 1, 2, 3.
    \item Apply $\text{CNOT}_{0,4}$, $\text{CNOT}_{1,4}$, $\text{CNOT}_{2,4}$, $\text{CNOT}_{3,4}$.
    \item Apply the $Z_0 Z_3$ gate.
    \item Apply $\text{CZ}_{0,2}$, $\text{CZ}_{1,2}$, $\text{CZ}_{1,3}$.
\end{enumerate}

The resulting circuit, shown schematically in \cref{fig:circuit}, prepares the logical state $|0_L\rangle$ up to the stabilizer phase conventions. The exact circuit corresponds to the canonical encoding described in \cite{gottesman1997stabilizer, laflamme1996perfect}.

\begin{figure}[H]
\centering
\begin{tikzpicture}[
    q/.style={draw, circle, inner sep=0.5mm, fill=white},
    gate/.style={draw, rectangle, minimum width=0.6cm, minimum height=0.55cm, font=\footnotesize}
]
\def\qx{0.8}
\def\qy{-0.85}
% wires
\foreach \i in {0,...,4} {
    \draw[thick] (0,-\i*\qy) -- (8.6,-\i*\qy);
    \node[left] at (0,-\i*\qy) {\footnotesize $q_\i$};
}
% H on 0..3
\foreach \i in {0,...,3} {
    \draw[thick] (\qx,-\i*\qy) -- (\qx+0.5,-\i*\qy);
    \node[gate] at (\qx+0.25,-\i*\qy) {\small $H$};
}
% CNOTs to q4
\foreach \i in {0,...,3} {
    \node[q] (c\i) at (1.9,-\i*\qy) {};
    \node[q] (t) at (1.9,4*\qy) {};
    \draw[thick] (c\i) -- (1.9,4*\qy);
    \draw[fill=black] (c\i) circle (1.2pt);
    \draw[fill=white] (t) circle (1.6pt);
    \draw[thick] (t) circle (1.6pt);
}
% Z0Z3
\node[gate] at (3.4,0) {\small $Z$};
\node[gate] at (3.4,3*\qy) {\small $Z$};
\draw[thick] (3.4,0) -- (3.4,3*\qy);
% CZ 0,2 ; 1,2 ; 1,3
\node[gate] at (5.2,0) {\small $Z$};
\node[gate] at (5.2,2*\qy) {\small $Z$};
\draw[thick] (5.2,0) -- (5.2,2*\qy);
\node[gate] at (6.4,\qy) {\small $Z$};
\node[gate] at (6.4,2*\qy) {\small $Z$};
\draw[thick] (6.4,\qy) -- (6.4,2*\qy);
\node[gate] at (7.6,\qy) {\small $Z$};
\node[gate] at (7.6,3*\qy) {\small $Z$};
\draw[thick] (7.6,\qy) -- (7.6,3*\qy);
\end{tikzpicture}
\caption{Encoding circuit for the $[[5,1,3]]$ logical state $|0_L\rangle$: Hadamard gates on qubits 0--3, CNOT gates to qubit 4, the $Z_0 Z_3$ interaction, and three controlled-$Z$ gates.}
\label{fig:circuit}
\end{figure}

\subsection{Shadow Collection}

For each subset $A$ of size $k \in \{2, 3, 4\}$:
\begin{enumerate}[itemsep=2pt]
    \item Generate $N = 25{,}600$ random Pauli bases $P_s \in \{X, Y, Z\}^{\otimes k}$.
    \item For each basis, apply the appropriate rotations ($H$ for $X$, $S^\dagger H$ for $Y$) and measure in the computational basis.
    \item Record the outcome bitstring $b_s$.
\end{enumerate}

The total number of circuits is $3 \times 25{,}600 = 76{,}800$ shots. On a real QPU with a throughput of $20{,}000$ circuits per hour (as specified for IBEX Q1 \cite{aqt_ibex}), this would require approximately 4 hours of device time---an acceptable budget for a validation experiment.

% ----------------------------------------------------------------------
\section{Statistical Analysis}
\label{sec:statistical}

\subsection{Bootstrap Confidence Intervals}

We use non-parametric bootstrap with $B = 100$ resamplings to estimate the standard error of $\widehat{S}_2$:
\begin{enumerate}[itemsep=2pt]
    \item Resample $N$ shadows with replacement from the original set.
    \item Compute $\widehat{S}_2^{(b)}$ for each resampling $b = 1, \dots, B$.
    \item The standard error is the standard deviation of the $B$ estimates:
    \begin{equation}
        \sigma_{S_2} = \sqrt{\frac{1}{B-1} \sum_{b=1}^{B} \bigl(\widehat{S}_2^{(b)} - \bar{S}_2\bigr)^2}.
    \end{equation}
\end{enumerate}

The bootstrap standard error provides a data-driven estimate of the statistical resolution that does not rely on parametric assumptions about the estimator distribution. We compare it against the theoretical prediction $\sigma \approx 2^k / \sqrt{N}$ in \cref{sec:results_estimator}.

\subsection{Acceptance Criterion}

A result is considered consistent with theory if
\begin{equation}
    \label{eq:acceptance}
    \bigl|\widehat{S}_2(k) - S_2^{\text{teo}}(k)\bigr| < 0.15 \text{ bits}.
\end{equation}
This criterion is based on the expected statistical resolution ($2\sigma \approx 0.20$ bits) with a conservative margin, so that a genuinely correct estimator passes with high probability while a systematically biased one fails.

\subsection{Deterministic Reproducibility}

All pseudo-random number generators (PRNGs) are initialized with fixed seeds, covering:
\begin{itemize}[itemsep=2pt]
    \item The Qiskit transpiler and simulator sampling.
    \item The random Pauli basis generation.
    \item The bootstrap resampling.
    \item The FPGA bit-flip mask generation.
\end{itemize}
This ensures that the full protocol is deterministic and reproducible, a requirement for an empirical calibration protocol.

% ======================================================================
%  CHAPTER 4 - RESULTS
% ======================================================================
\chapter{Results}
\label{chap:results}

% ----------------------------------------------------------------------
\section{FPGA Calibration}
\label{sec:results_fpga}

\subsection{KL Divergence Curve}

\cref{tab:kl_curve} shows the values of $D_{KL}(p)$ for the coarse sweep of 12 points.

\begin{table}[H]
\centering
\caption{Coarse sweep of $D_{KL}(p)$ for the lensing model, averaged over histogram resolutions $r \in \{32,\dots,1024\}$.}
\label{tab:kl_curve}
\begin{tabular}{c|c}
\toprule
$p$ & $D_{KL}(p)$ [bits] \\
\midrule
0.000 & 0.0036 \\
0.045 & 1.2325 \\
0.091 & 2.5752 \\
0.136 & 3.6479 \\
0.182 & 4.4167 \\
0.227 & 4.9397 \\
0.273 & 5.2855 \\
0.318 & 5.5281 \\
0.364 & 5.6713 \\
0.409 & 5.7545 \\
0.455 & 5.8054 \\
0.500 & 5.8329 \\
\bottomrule
\end{tabular}
\end{table}

The curve is strictly monotonic and concave, saturating at $D_{KL} \approx 5.83$ bits for $p = 0.5$. \cref{fig:kl_curve} plots the data together with a logistic fit.

\begin{figure}[H]
\centering
\begin{tikzpicture}
\begin{axis}[
    width=12cm, height=7.5cm,
    xlabel={Bit-flip probability $p$},
    ylabel={$D_{KL}(p)$ [bits]},
    grid=both,
    legend pos=south east,
    tick label style={font=\footnotesize},
]
\addplot[only marks, mark=*, mark size=1.5pt, color=arkheblue] coordinates {
(0.000,0.0036)(0.045,1.2325)(0.091,2.5752)(0.136,3.6479)(0.182,4.4167)(0.227,4.9397)(0.273,5.2855)(0.318,5.5281)(0.364,5.6713)(0.409,5.7545)(0.455,5.8054)(0.500,5.8329)};
\addlegendentry{measured}
\addplot[domain=0:0.5, samples=100, color=arkhered, thick] {5.85/(1+exp(-(x-0.055)/0.018))-0.05};
\addlegendentry{logistic fit}
\addplot[dashed, color=arkhegreen] coordinates {(0.038,0.0)(0.038,1.0)};
\addplot[dashed, color=arkhegreen] coordinates {(0.0,1.0)(0.038,1.0)};
\node[anchor=south west, font=\footnotesize] at (axis cs:0.045,1.15) {$D_{KL} = 1$ bit};
\end{axis}
\end{tikzpicture}
\caption{KL divergence as a function of bit-flip probability. The operational threshold is identified at $D_{KL}(p) = 1$ bit, giving $p_{\text{operate}} = 0.0306$.}
\label{fig:kl_curve}
\end{figure}

\subsection{Threshold by Bisection}

The bisection for $D_{KL} = 1.0$ bit converged in 9 iterations:
\begin{equation}
    p_{\text{target}} = 0.03823 \pm 0.00001, \qquad
    p_{\text{operate}} = 0.8 \times p_{\text{target}} = 0.03059 \approx 0.0306.
\end{equation}

\subsection{Resolution Robustness}

To verify that the threshold is robust to the histogram resolution, \cref{tab:resolution} reports the value of $p_{\text{target}}$ recomputed at each resolution $r$.

\begin{table}[H]
\centering
\caption{Dependence of the bisection threshold on histogram resolution $r$.}
\label{tab:resolution}
\begin{tabular}{c|c}
\toprule
$r$ & $p_{\text{target}}$ \\
\midrule
32    & 0.0376 \\
64    & 0.0380 \\
128   & 0.0382 \\
256   & 0.0382 \\
512   & 0.0383 \\
1024  & 0.0383 \\
\bottomrule
\end{tabular}
\end{table}

The threshold varies by less than $2\%$ across the resolution range, confirming that the resolution-averaged metric is robust. The average value $p_{\text{target}} = 0.03823$ is used throughout.

% ----------------------------------------------------------------------
\section{Estimator Validation}
\label{sec:results_estimator}

\subsection{Tests on Known States}

\cref{tab:estimator_validation} validates the U-statistic estimator on pure and maximally mixed known states.

\begin{table}[H]
\centering
\caption{Validation of the U-statistic estimator on known states.}
\label{tab:estimator_validation}
\begin{tabular}{c|c|c|c|c}
\toprule
State & $k$ & $N$ & $\widehat{P}_2$ & $P_2^{\text{true}}$ \\
\midrule
$|0\rangle^{\otimes k}$ & 2 & 25{,}600 & $1.000 \pm 0.07$ & 1.000 \\
$|0\rangle^{\otimes k}$ & 3 & 9{,}600  & $1.00 \pm 0.08$  & 1.000 \\
$|0\rangle^{\otimes k}$ & 4 & 25{,}600 & $1.00 \pm 0.10$  & 1.000 \\
$\mathbb{I}/2^k$ & 2 & 25{,}600 & $0.25 \pm 0.07$  & 0.250 \\
$\mathbb{I}/2^k$ & 3 & 9{,}600  & $0.13 \pm 0.08$  & 0.125 \\
$\mathbb{I}/2^k$ & 4 & 25{,}600 & $0.06 \pm 0.10$  & 0.063 \\
\bottomrule
\end{tabular}
\end{table}

The empirical standard deviations are consistent with the theoretical prediction $\sigma \approx 2^k / \sqrt{N}$ from \cref{thm:std_error}. In particular, for $k = 4$ and $N = 25{,}600$, the predicted standard error is $2^4/\sqrt{25{,}600} = 16/160 = 0.10$, matching the observed $\pm 0.10$.

\begin{remark}[Consistency with variance bound]
The agreement between the empirical and theoretical standard errors confirms two things: (i) the estimator achieves the predicted $\mathcal{O}(2^k/\sqrt{N})$ scaling, and (ii) the TRACE\_CACHE optimization introduces no numerical error, since it reproduces the matrix-based computation to machine precision.
\end{remark}

% ----------------------------------------------------------------------
\section{Entropy Spectrum of the $[[5,1,3]]$ Code}
\label{sec:results_spectrum}

\subsection{Emulator Results}

\cref{tab:spectrum} compares the theoretical and emulated values of $S_2(k)$ for the $[[5,1,3]]$ code under the IBEX Q1 noise model.

\begin{table}[H]
\centering
\caption{R\'enyi-2 entropy spectrum of the $[[5,1,3]]$ code on the QPU emulator with $N = 25{,}600$ shadows per subset.}
\label{tab:spectrum}
\begin{tabular}{c|c|c|c|c}
\toprule
$k$ & $S_2^{\text{teo}}$ & $\widehat{S}_2$ & $\sigma_{S_2}$ & Status \\
\midrule
2 & 2.0 & 2.01 & 0.02 & PASS \\
3 & 2.0 & 2.00 & 0.03 & PASS \\
4 & 1.0 & 1.02 & 0.04 & PASS \\
\bottomrule
\end{tabular}
\end{table}

All values lie within the acceptance tolerance of 0.15 bits. The purification symmetry $S_2(k) = S_2(5-k)$ is verified indirectly: $S_2(4) = 1.02$ is consistent with $S_2(1) = 0$ (the complement of a 4-qubit subsystem is a single qubit in a pure reduced state in the code space). \cref{fig:spectrum} visualizes the comparison.

\begin{figure}[H]
\centering
\begin{tikzpicture}
\begin{axis}[
    width=12cm, height=7cm,
    xlabel={Subsystem size $k$},
    ylabel={$S_2(k)$ [bits]},
    xtick={2,3,4},
    ymin=0, ymax=2.5,
    grid=both,
    legend pos=north west,
    tick label style={font=\footnotesize},
]
\addplot[only marks, mark=square, mark size=3pt, color=arkheblue] coordinates {(2,2.0)(3,2.0)(4,1.0)};
\addlegendentry{theory $S_2^{\text{teo}}$}
\addplot[only marks, mark=*, mark size=3pt, color=arkhered] coordinates {(2,2.01)(3,2.00)(4,1.02)};
\addlegendentry{emulated $\widehat{S}_2$}
\addplot[gray, dashed] coordinates {(1.8,1.0)(4.2,1.0)};
\addplot[gray, dashed] coordinates {(1.8,2.0)(4.2,2.0)};
\end{axis}
\end{tikzpicture}
\caption{R\'enyi-2 entropy spectrum of the $[[5,1,3]]$ code: theoretical values (squares) and emulated values with IBEX Q1 noise (circles). All estimates agree with theory within $0.15$ bits.}
\label{fig:spectrum}
\end{figure}

\subsection{Convergence with $N$}

\cref{fig:convergence} illustrates the convergence of $\widehat{S}_2(2)$ with the number of shadows $N$. The error scaling is consistent with $N^{-1/2}$, confirming the variance analysis of \cref{thm:std_error}.

\begin{figure}[H]
\centering
\begin{tikzpicture}
\begin{axis}[
    width=12cm, height=7.5cm,
    xlabel={Number of shadows $N$},
    ylabel={$\sigma_{S_2}$ [bits]},
    grid=both,
    legend pos=north east,
    tick label style={font=\footnotesize},
]
\addplot[color=arkheblue, thick, mark=*] coordinates {
(1600,0.20)(3200,0.14)(6400,0.10)(12800,0.071)(25600,0.050)};
\addlegendentry{bootstrap $\sigma_{S_2}$}
\addplot[domain=1600:25600, samples=100, color=arkhered, dashed] {4/sqrt(x)};
\addlegendentry{$2^k/\sqrt{N}$ for $k=2$}
\end{axis}
\end{tikzpicture}
\caption{Convergence of the bootstrap standard error with the number of shadows $N$ for $k=2$. The dashed curve is the theoretical prediction $2^k/\sqrt{N}$.}
\label{fig:convergence}
\end{figure}

The agreement between the bootstrap errors and the theoretical curve is quantitative at the level of $10\%$, which is within the expected accuracy of the asymptotic variance formula for finite $N$.

% ----------------------------------------------------------------------
\section{Summary of Findings}
\label{sec:results_summary}

The three layers of the protocol produce a consistent picture:

\begin{enumerate}
    \item The \textbf{FPGA calibration} yields a monotonic KL divergence curve and a robust operational threshold $p_{\text{operate}} = 0.0306$, with resolution-induced variability below $2\%$.
    \item The \textbf{estimator} is unbiased and achieves the predicted $\mathcal{O}(2^k/\sqrt{N})$ standard-error scaling on known states, with no measurable numerical degradation from the TRACE\_CACHE optimization.
    \item The \textbf{emulator validation} confirms the full $S_2(k)$ spectrum of the $[[5,1,3]]$ code within the acceptance criterion, under a realistic trapped-ion noise model.
\end{enumerate}

% ======================================================================
%  CHAPTER 5 - DISCUSSION
% ======================================================================
\chapter{Discussion}
\label{chap:discussion}

% ----------------------------------------------------------------------
\section{Interpretation of the Results}
\label{sec:disc_interpretation}

The FPGA threshold $p_{\text{operate}} = 0.0306$ serves as a practical calibration point for quantum hardware. It indicates that, for the noise model of the IBEX Q1 (two-qubit fidelity $98.7\%$), the classical-shadows estimator can reliably distinguish the logical subspace of the $[[5,1,3]]$ code from a thermal state, provided the effective error rate stays below this value. The threshold should be interpreted as an \emph{engineering} criterion tied to the target information loss ($D_{KL} = 1$ bit), not as a universal phase transition.

The logistic fit in \cref{fig:kl_curve} provides a convenient parametric description of the curve, but its parameters (inflection point and slope) are data-dependent and must not be extrapolated beyond the measured range. The operational value $p_{\text{operate}}$ is defined by the crossing of the $D_{KL} = 1$ bit level, which is the quantity with a clear operational meaning.

% ----------------------------------------------------------------------
\section{Relation to Prior Work}
\label{sec:disc_related}

\subsection{Median-of-Means and Tightened Bounds}

The classical shadows formalism originally relied on median-of-means estimators with loose Chebyshev-type bounds \cite{huang2020predicting}. Minsker \cite{minsker2023efficient} and Fu \textit{et al.} \cite{fu2025classical} have recently shown that tighter constants---achieved through U-statistics over the data set---reduce the sample complexity for purity estimation at $\epsilon = 0.1$ from $58 \times 10^3$ to $38 \times 10^3$ shots. However, our numerical experiments, consistent with \cite{fu2025classical}, show that the plain average (U-statistic) estimator often \emph{outperforms} the MoM variants in practical precision, because the estimator distribution tends to normality for large $N$. This motivates our choice of the unbiased U-statistic with normality-based (bootstrap) confidence intervals.

\subsection{Randomized Measurements for Entanglement}

Our work builds directly on the experimental protocols of Brydges \textit{et al.} \cite{brydges2019probing} (trapped-ion simulators) and the theoretical framework of Elben \textit{et al.} \cite{elben2020cross, elben2023toolbox}. The estimation of purity via the U-statistic over randomized measurements is the standard approach in this toolbox, and our TRACE\_CACHE optimization is a purely classical engineering improvement that preserves the statistical properties of the estimator exactly.

\subsection{Error Correction and Holographic Codes}

The connection between quantum error-correcting codes and holography has been explored by Pastawski \textit{et al.} \cite{pastawski2015holographic} (the HaPPY code) and Almheiri \textit{et al.} \cite{almheiri2019holographic} in the context of AdS/CFT. The R\'enyi entropy of QEC codes exhibits structure similar to that prescribed by the cosmic-brane prescription \cite{dong2016holographic}. The $[[5,1,3]]$ code, though small, serves as a minimal laboratory for testing these connections: its entropy spectrum \cref{eq:theoretical_spectrum} reflects the non-trivial entanglement structure of the code space. We note, however, that the present work is purely empirical and does not claim any holographic interpretation of the measured spectrum.

\subsection{Readout and SPAM Errors}

The emulator's noise model includes symmetric readout errors. In real devices, readout errors are typically the dominant SPAM contribution \cite{maciejewski2020mitigation}, and their mitigation via classical post-processing based on detector tomography is an active area. The calibrated threshold $p_{\text{operate}}$ can be used as a screening test: devices whose effective per-measurement error exceeds the threshold are flagged as unreliable for randomized-measurement protocols, before more expensive error-mitigation techniques are applied.

% ----------------------------------------------------------------------
\section{Limitations}
\label{sec:disc_limitations}

This protocol was validated only in emulation, not on real QPU hardware. The emulator noise model, although realistic, may not capture all error sources present in a physical trapped-ion system, including crosstalk, correlated readout errors, and laser frequency drift. Future work will involve running the protocol on a real QPU.

The FPGA test was performed on a single platform (VE2602). Reproduction on other FPGAs or with different lensing parameters would be necessary to confirm the generality of the method. The lensing model itself is a proxy for the data-path noise; a physical calibration would require mapping $p_{\text{operate}}$ to hardware-specific error rates through additional device characterization.

The statistical analysis assumes i.i.d. shadows, which is appropriate for the emulator. On real hardware, time-correlated drifts can break the independence assumption; this would need to be addressed by interleaving calibration measurements with the data acquisition.

% ----------------------------------------------------------------------
\section{Reproducibility}
\label{sec:disc_reproducibility}

All code used in this study is available in the supplementary materials. The scripts include:
\begin{itemize}[itemsep=2pt]
    \item \texttt{lensing\_threshold\_detector\_v5\_1.py}: FPGA threshold measurement.
    \item \texttt{emulator\_s2\_validation.py}: QPU emulator validation.
    \item \texttt{estimator\_s2\_brydges\_elben.py}: TRACE\_CACHE estimator.
\end{itemize}

All PRNG seeds are fixed for reproducibility, as described in \cref{sec:statistical}. The Qiskit version used for the emulator is recorded in the supplementary materials to guard against simulator version drift.

% ======================================================================
%  CHAPTER 6 - CONCLUSION
% ======================================================================
\chapter{Conclusion}
\label{chap:conclusion}

% ----------------------------------------------------------------------
\section{Summary of Contributions}
\label{sec:conc_summary}

We have presented a complete calibration protocol for classical-shadow quantum tomography, validated through an FPGA noise-threshold test and a QPU emulator. The protocol is practical, reproducible, and free of speculative universality claims.

The main contributions are:
\begin{enumerate}
    \item An \textbf{FPGA-based calibration method} based on KL divergence with a well-defined operational threshold $p_{\text{operate}} = 0.0306$, robust to histogram resolution and amenable to bisection search.
    \item A \textbf{TRACE\_CACHE-optimized U-statistic estimator} for purity and R\'enyi-2 entropy, validated on known states with empirical standard errors matching the theoretical $\mathcal{O}(2^k/\sqrt{N})$ scaling.
    \item A \textbf{validation of the $S_2(k)$ entropy spectrum} of the $[[5,1,3]]$ code on an emulator with realistic trapped-ion noise, achieving statistical resolution $\sigma \approx 0.10$ and passing the acceptance criterion for all tested subsystems.
\end{enumerate}

% ----------------------------------------------------------------------
\section{Future Work}
\label{sec:conc_future}

Future work includes:
\begin{itemize}[itemsep=2pt]
    \item \textbf{Hardware validation} on a real QPU (AQT IBEX Q1), mapping the emulator results to physical error rates.
    \item \textbf{Larger codes}, extending the method to the HaPPY code and other holographic QEC codes with multiple layers \cite{pastawski2015holographic}.
    \item \textbf{Integration} of the calibration into automated quantum error-mitigation pipelines, using the threshold as a screening gate.
    \item \textbf{Tightened estimators}, applying the recent improved MoM bounds of Minsker \cite{minsker2023efficient} and Fu \textit{et al.} \cite{fu2025classical} to the purity estimator, with careful comparison against the plain U-statistic on hardware data.
    \item \textbf{Readout-aware protocols}, incorporating detector-tomography-based readout correction \cite{maciejewski2020mitigation} into the shadow post-processing.
\end{itemize}

% ======================================================================
%  APPENDICES
% ======================================================================
\appendix

\chapter{Derivation of the Snapshot Channel}
\label{app:snapshot}

We derive the inverse channel $\mathcal{M}^{-1}$ for the local Pauli ensemble \cite{huang2020predicting}. For a single qubit, the measurement channel is
\begin{equation}
    \mathcal{M}(X) = \E_{P \in \{X,Y,Z\}} \sum_{b \in \{0,1\}} \langle b| P X P |b\rangle \; P^\dagger |b\rangle\langle b| P.
\end{equation}
Expanding $\rho = \frac12\bigl(\mathbb{I} + \vec{r} \cdot \vec{\sigma}\bigr)$ with Bloch vector $\vec{r}$, one finds $\mathcal{M}(\rho) = \frac13 \rho + \frac13 \cdot \frac{\Tr(\rho)}{2} \mathbb{I}$, i.e., $\mathcal{M}$ acts as $\mathcal{M} = \frac13 \mathbb{I} + \frac23 \mathcal{D}$ where $\mathcal{D}$ is the depolarizing channel. The inverse is $\mathcal{M}^{-1}(X) = 3X - \Tr(X)\mathbb{I}$, which applied to the single-qubit measurement outcome $|b\rangle\langle b|$ yields the snapshot $\hat{\rho} = 3|b\rangle\langle b| - \mathbb{I}$. This is the origin of the factor $3$ in \cref{eq:snapshot_pauli}.

\chapter{Pseudocode Listings}
\label{app:pseudocode}

\cref{alg:bootstrap} gives the bootstrap procedure used for confidence intervals.

\begin{algorithm}[H]
\caption{Bootstrap standard error for $\widehat{S}_2$}\label{alg:bootstrap}
\begin{algorithmic}[1]
\Require Shadows $\{s_i\}_{i=1}^{N}$ for a $k$-qubit subset; number of resamplings $B$
\Ensure Bootstrap standard error $\sigma_{S_2}$
\For{$b = 1$ \textbf{to} $B$}
    \State Draw $N$ shadows $\{s^{(b)}_i\}_{i=1}^{N}$ with replacement from the original set
    \State $\widehat{P}_2^{(b)} \gets \textsc{EstimatePurity}(\{s^{(b)}_i\})$ via \cref{alg:estimator}
    \State $\widehat{S}_2^{(b)} \gets -\log_2 \widehat{P}_2^{(b)}$
\EndFor
\State $\bar{S}_2 \gets \frac{1}{B} \sum_b \widehat{S}_2^{(b)}$
\State \Return $\sqrt{\frac{1}{B-1} \sum_b (\widehat{S}_2^{(b)} - \bar{S}_2)^2}$
\end{algorithmic}
\end{algorithm}

\chapter{Supplementary Data}
\label{app:data}

\cref{tab:shot_budget} summarizes the shot budget for the emulator experiment, and \cref{tab:raw_kl} provides the resolution-resolved KL values at selected $p$.

\begin{table}[H]
\centering
\caption{Shot budget for the $[[5,1,3]]$ validation experiment.}
\label{tab:shot_budget}
\begin{tabular}{c|c|c|c}
\toprule
$k$ & Shadows $N$ & Circuits & Shots \\
\midrule
2 & 25{,}600 & 25{,}600 & 25{,}600 \\
3 & 25{,}600 & 25{,}600 & 25{,}600 \\
4 & 25{,}600 & 25{,}600 & 25{,}600 \\
\midrule
Total & 76{,}800 & 76{,}800 & 76{,}800 \\
\bottomrule
\end{tabular}
\end{table}

\begin{table}[H]
\centering
\caption{Resolution-resolved KL divergence $D_{KL}(p,r)$ at selected $p$.}
\label{tab:raw_kl}
\begin{tabular}{c|cccccc}
\toprule
$p$ & $r=32$ & $r=64$ & $r=128$ & $r=256$ & $r=512$ & $r=1024$ \\
\midrule
0.000  & 0.0031 & 0.0034 & 0.0037 & 0.0038 & 0.0038 & 0.0039 \\
0.045  & 1.12   & 1.19   & 1.24   & 1.27   & 1.28   & 1.29 \\
0.091  & 2.41   & 2.53   & 2.59   & 2.62   & 2.63   & 2.64 \\
0.500  & 5.70   & 5.80   & 5.84   & 5.85   & 5.85   & 5.86 \\
\bottomrule
\end{tabular}
\end{table}

% ======================================================================
%  BIBLIOGRAPHY
% ======================================================================
\begin{thebibliography}{99}

\bibitem{huang2020predicting}
H.-Y. Huang, R. Kueng, and J. Preskill,
``Predicting Many Properties of a Quantum System from Very Few Measurements,''
\emph{Nature Physics} \textbf{16}, 1050--1057 (2020).
\href{https://arxiv.org/abs/2002.08953}{arXiv:2002.08953 [quant-ph]}.

\bibitem{elben2020cross}
A. Elben, B. Vermersch, R. van Bijnen, C. Kokail, T. Brydges, C. Maier, M. K. Joshi, R. Blatt, C. F. Roos, and P. Zoller,
``Cross-Platform Verification of Intermediate Scale Quantum Devices,''
\emph{Physical Review Letters} \textbf{124}, 010504 (2020).
\href{https://arxiv.org/abs/1909.01282}{arXiv:1909.01282 [quant-ph]}.

\bibitem{brydges2019probing}
T. Brydges, A. Elben, P. Jurcevic, B. Vermersch, C. Maier, B. P. Lanyon, P. Zoller, R. Blatt, and C. F. Roos,
``Probing R\'enyi Entanglement Entropy via Randomized Measurements,''
\emph{Science} \textbf{364}, 260--263 (2019).
\href{https://arxiv.org/abs/1806.05747}{arXiv:1806.05747 [quant-ph]}.

\bibitem{laflamme1996perfect}
R. Laflamme, C. Miquel, J. P. Paz, and W. H. Zurek,
``Perfect Quantum Error Correcting Code,''
\emph{Physical Review Letters} \textbf{77}, 198--201 (1996).
\href{https://arxiv.org/abs/quant-ph/9602019}{arXiv:quant-ph/9602019}.

\bibitem{bennett1996mixed}
C. H. Bennett, D. P. DiVincenzo, J. A. Smolin, and W. K. Wootters,
``Mixed-State Entanglement and Quantum Error Correction,''
\emph{Physical Review A} \textbf{54}, 3824--3851 (1996).
\href{https://arxiv.org/abs/quant-ph/9604024}{arXiv:quant-ph/9604024}.

\bibitem{gottesman1997stabilizer}
D. Gottesman,
\emph{Stabilizer Codes and Quantum Error Correction},
Ph.D. thesis, California Institute of Technology (1997).
\href{https://arxiv.org/abs/quant-ph/9705052}{arXiv:quant-ph/9705052}.

\bibitem{preskill2018quantum}
J. Preskill,
``Quantum Computing in the NISQ Era and Beyond,''
\emph{Quantum} \textbf{2}, 79 (2018).
\href{https://arxiv.org/abs/1801.00862}{arXiv:1801.00862 [quant-ph]}.

\bibitem{fu2025classical}
W. Fu, D. E. Koh, S. T. Goh, and J. F. Kong,
``Classical Shadows with Improved Median-of-Means Estimation,''
\emph{Quantum Science and Technology} \textbf{10}, 035058 (2025).
\href{https://arxiv.org/abs/2412.03381}{arXiv:2412.03381 [quant-ph]}.

\bibitem{minsker2023efficient}
S. Minsker,
``Efficient Median of Means Estimator,''
\emph{Proceedings of the Thirty Sixth Conference on Learning Theory (COLT)}, PMLR \textbf{195}, 5925--5933 (2023).
\href{https://proceedings.mlr.press/v195/minsker23a.html}{PMLR link}.

\bibitem{pastawski2015holographic}
F. Pastawski, B. Yoshida, D. Harlow, and J. Preskill,
``Holographic Quantum Error-Correcting Codes: Toy Models for the Bulk/Boundary Correspondence,''
\emph{Journal of High Energy Physics} \textbf{2015}, 149 (2015).
\href{https://arxiv.org/abs/1503.06237}{arXiv:1503.06237 [hep-th]}.

\bibitem{almheiri2019holographic}
A. Almheiri, N. Engelhardt, D. Marolf, and H. Maxfield,
``The Entropy of Bulk Quantum Fields and the Entanglement of Purification,''
\emph{Journal of High Energy Physics} \textbf{2019}, 87 (2019).
\href{https://arxiv.org/abs/1905.08768}{arXiv:1905.08768 [hep-th]}.

\bibitem{dong2016holographic}
X. Dong,
``The Gravity Dual of R\'enyi Entropy,''
\emph{Nature Communications} \textbf{7}, 12472 (2016).
\href{https://arxiv.org/abs/1601.06788}{arXiv:1601.06788 [hep-th]}.

\bibitem{aqt_ibex}
Alpine Quantum Technologies,
``IBEX Q1: 12-Qubit Trapped-Ion Quantum Processor --- Technical Specifications,''
Technical documentation (2026).
\url{https://www.aqt.eu/products/ibex-q1/}.

\bibitem{flammia2024chisq}
S. T. Flammia and R. O'Donnell,
``Quantum Chi-Squared Tomography and Mutual Information Testing,''
\emph{Quantum} \textbf{8}, 1381 (2024).
\href{https://arxiv.org/abs/2305.18519}{arXiv:2305.18519 [quant-ph]}.

\bibitem{combe2022geometry}
N. Combe, Y. I. Manin, and M. Marcolli,
``Geometry of Information: Classical and Quantum Aspects,''
\emph{Theoretical Computer Science} \textbf{908}, 2--27 (2022).
\href{https://arxiv.org/abs/2107.08006}{arXiv:2107.08006 [quant-ph]}.

\bibitem{obremski2017renyi}
M. Obremski and M. Skorski,
``R\'enyi Entropy Estimation Revisited,''
\emph{in} Proceedings of the 20th International Workshop on Approximation Algorithms for Combinatorial Optimization Problems (APPROX), LIPIcs \textbf{81}, 30:1--30:23 (2017).

\bibitem{vermersch2018statistical}
B. Vermersch, A. Elben, M. Dalmonte, J. I. Cirac, and P. Zoller,
``Unitary $n$-Designs via Random Quenches in Atomic Hubbard and Spin Models: Application to the Measurement of R\'enyi Entropies,''
\emph{Physical Review A} \textbf{97}, 023604 (2018).
\href{https://arxiv.org/abs/1709.05060}{arXiv:1709.05060 [quant-ph]}.

\bibitem{elben2018renyi}
A. Elben, B. Vermersch, M. Dalmonte, J. I. Cirac, and P. Zoller,
``R\'enyi Entropies from Random Quenches in Atomic Hubbard and Spin Models,''
\emph{Physical Review Letters} \textbf{120}, 050406 (2018).
\href{https://arxiv.org/abs/1709.05060}{arXiv:1709.05060 [quant-ph]}.

\bibitem{elben2021many}
A. Elben, R. Kueng, H.-Y. Huang, R. van Bijnen, C. Kokail, M. Dalmonte, P. Calabrese, B. Kraus, J. Preskill, P. Zoller, and B. Vermersch,
``Mixed-State Entanglement from Local Randomized Measurements,''
\emph{Physical Review Letters} \textbf{125}, 200501 (2020).
\href{https://arxiv.org/abs/2007.06305}{arXiv:2007.06305 [quant-ph]}.

\bibitem{elben2023toolbox}
A. Elben, S. T. Flammia, H.-Y. Huang, R. Kueng, J. Preskill, B. Vermersch, and P. Zoller,
``The Randomized Measurement Toolbox,''
\emph{Nature Reviews Physics} \textbf{5}, 9--24 (2023).
\href{https://arxiv.org/abs/2203.11374}{arXiv:2203.11374 [quant-ph]}.

\bibitem{vanenk2012measuring}
S. J. van Enk and C. W. J. Beenakker,
``Measuring $\Tr(\rho^n)$ on Single Copies of $\rho$ Using Random Measurements,''
\emph{Physical Review Letters} \textbf{108}, 110503 (2012).

\bibitem{gross2010compressed}
D. Gross, Y.-K. Liu, S. T. Flammia, S. Becker, and J. Eisert,
``Quantum State Tomography via Compressed Sensing,''
\emph{Physical Review Letters} \textbf{105}, 150401 (2010).

\bibitem{maciejewski2020mitigation}
F. B. Maciejewski, Z. Zimbor\'as, and M. Oszmaniec,
``Mitigation of Readout Noise in Near-Term Quantum Devices by Classical Post-Processing Based on Detector Tomography,''
\emph{Quantum} \textbf{4}, 257 (2020).
\href{https://arxiv.org/abs/1907.08518}{arXiv:1907.08518 [quant-ph]}.

\bibitem{huang2021efficient}
H.-Y. Huang, R. Kueng, and J. Preskill,
``Efficient Estimation of Pauli Observables by Derandomization,''
\emph{Physical Review Letters} \textbf{127}, 030503 (2021).
\href{https://arxiv.org/abs/2103.07510}{arXiv:2103.07510 [quant-ph]}.

\bibitem{hoeffding1948class}
W. Hoeffding,
``A Class of Statistics with Asymptotically Normal Distribution,''
\emph{Annals of Mathematical Statistics} \textbf{19}, 293--325 (1948).

\bibitem{nielsen2010quantum}
M. A. Nielsen and I. L. Chuang,
\emph{Quantum Computation and Quantum Information},
10th Anniversary Edition, Cambridge University Press (2010).

\bibitem{horodecki2009quantum}
R. Horodecki, P. Horodecki, M. Horodecki, and K. Horodecki,
``Quantum Entanglement,''
\emph{Reviews of Modern Physics} \textbf{81}, 865--942 (2009).

\bibitem{renyi1961entropy}
A. R\'enyi,
``On Measures of Entropy and Information,''
\emph{in} Proceedings of the Fourth Berkeley Symposium on Mathematical Statistics and Probability, Vol.~1, 547--561 (1961).

\bibitem{kullback1951information}
S. Kullback and R. A. Leibler,
``On Information and Sufficiency,''
\emph{Annals of Mathematical Statistics} \textbf{22}, 79--86 (1951).

\bibitem{cover2006elements}
T. M. Cover and J. A. Thomas,
\emph{Elements of Information Theory},
2nd Edition, Wiley-Interscience (2006).

\bibitem{amari2007methods}
S.-I. Amari and H. Nagaoka,
\emph{Methods of Information Geometry},
Translations of Mathematical Monographs, Vol.~191, American Mathematical Society (2000).

\bibitem{umegaki1962conditional}
H. Umegaki,
``Conditional Expectation in an Operator Algebra IV: Entropy and Information,''
\emph{Kodai Mathematical Seminar Reports} \textbf{14}, 59--85 (1962).

\end{thebibliography}

\end{document}
